% !PRINTED 280
\weberpage{280}
We can therefore apply the same procedure as above if we add the further
restriction that only \gesperrt{symmetric} functions of $A$ and $B$, that
is, rational functions of $AB$, $A+B$, are used, because \gesperrt{only} under this
assumption do all the others follow from one of the values of such a
% [sic] The source cites (17); the relevant permutation table is (21), while (17) is the preceding list of transformed arguments.
function by the permutations (17). The following examples are easily
computed:
\weq{22}{
\begin{array}{l}
n = 15, \quad AB+1 = 0, \\
n = 35, \quad 2(A+B)-AB = 5 \\
n = 39, \quad 2(A+B)-AB = 3.
\end{array}
}

The computation presents no insurmountable difficulties even in still
more complicated cases. Thus in the memoir Acta mathematica Bd. II,
S. 359 I have communicated the following formula for the case $n = 105$:
\[
\begin{array}{l}
A = \sum \dfrac{f(3\omega)f(5\omega)f(7\omega)f(105\omega)}{f(\omega)f(15\omega)f(21\omega)f(35\omega)}, \\[2ex]
B = \sum \dfrac{f(\omega)f(15\omega)f(21\omega)f(35\omega)}{f(3\omega)f(5\omega)f(7\omega)f(105\omega)}, \\[2ex]
A^2+B^2-4(A+B)^3+10AB(A+B) \\
\quad{}+4(A+B)^2+10AB+14(A+B)+5 = 0,
\end{array}
\]
if the sums $\Sigma$ in $A$ and $B$ refer to the three functions
$f, f_1, f_2$.

\begin{center}
§ 77.\quad \gesperrt{Geometric interpretation of the irrational modular equations as}\\
\gesperrt{modular correspondences.}
\end{center}

We make the content of the irrational forms of the modular equations
more intuitive by means of a geometric interpretation. Let us first
consider the case $n \equiv -1 \ (\mathrm{mod}\ 8)$ and put
\weq{1}{x = f(\omega), \quad y = f_1(\omega), \quad z = f_2(\omega),}
so that, if we regard $x, y, z$ as Cartesian coordinates of a point
$P$, a space curve is thereby represented, which we can also express
by the two equations:
\weq{2}{
\begin{array}{c}
x^8 - y^8 - z^8 = 0 \\
xyz = \sqrt{2}
\end{array}
}
and which we shall call the \gesperrt{fundamental curve}. This curve is
therefore of the 24th order. One can map it onto
%CONT>>
